% =====================================================================
%  CARNET D'ENTRAÎNEMENT — CHIMIE — FICHE 3 — THÈME 2
%  Composition : atomes et ions — NIVEAU MOYEN — CORRIGÉ
% =====================================================================
\documentclass[11pt,a4paper]{article}

\usepackage[utf8]{inputenc}
\usepackage[T1]{fontenc}
\usepackage{lmodern}
\usepackage[french]{babel}

\usepackage{amsmath,amssymb,mathtools}
\usepackage{xcolor}
\usepackage{colortbl}
\usepackage{tikz}
\usetikzlibrary{arrows.meta,positioning,calc}
\usepackage[most]{tcolorbox}
\usepackage{enumitem}
\usepackage{booktabs}
\usepackage{array}
\usepackage[a4paper,margin=2.1cm,headheight=26pt,headsep=9pt]{geometry}
\usepackage{fancyhdr}
\usepackage[hidelinks]{hyperref}

\definecolor{primary}{RGB}{124,78,140}
\definecolor{primarydark}{RGB}{74,44,92}
\definecolor{defcol}{RGB}{0,114,178}
\definecolor{thmcol}{RGB}{0,140,120}
\definecolor{excol}{RGB}{0,135,90}
\definecolor{attncol}{RGB}{200,40,55}

\newcommand{\nuc}[3]{\prescript{#1}{#2}{\mathrm{#3}}}
\newcommand{\pp}{p^{+}}
\newcommand{\ee}{e^{-}}
\newcommand{\nz}{n^{0}}
\newcommand{\rep}{\textbf{\color{excol}Réponse : }}

\pagestyle{fancy}
\fancyhf{}
\fancyhead[L]{\small\textcolor{primarydark}{\textbf{Fiche 3 — Thème 2}} : Composition des atomes et des ions}
\fancyhead[R]{\small\textcolor{primarydark}{\textsc{Moyen — Corrigé}}}
\fancyfoot[C]{\small\thepage}
\renewcommand{\headrulewidth}{0.8pt}
\renewcommand{\headrule}{\hbox to\headwidth{\color{primary}\leaders\hrule height \headrulewidth\hfill}}

\tcbset{
  boxbase/.style={enhanced,breakable,boxrule=0.8pt,arc=2pt,
     left=2.6mm,right=2.6mm,top=1.8mm,bottom=1.8mm,
     fonttitle=\bfseries,coltitle=white,toptitle=0.4mm,bottomtitle=0.4mm},
}
\newtcolorbox{bilan}[1][]{boxbase,colback=thmcol!7,colframe=thmcol,
  title={\textsc{À retenir}\ifx\relax#1\relax\relax\else\textmd{ : \itshape #1}\fi}}

\newcounter{exo}
\newcommand{\cor}[1][]{\stepcounter{exo}\par\medskip%
  \noindent{\color{primary}\rule[-2pt]{3.5pt}{15pt}}\ \textbf{\color{primarydark}\large Exercice \theexo}%
  \ \textmd{\normalsize\color{black!55}— corrigé}%
  \ifx\relax#1\relax\relax\else\ \textmd{\normalsize\color{black!65}\itshape (#1)}\fi\par\nopagebreak\smallskip}

\setlist{leftmargin=1.7em,topsep=2pt,itemsep=3pt}

\begin{document}

\begin{tikzpicture}
\node[rectangle,rounded corners=6pt,fill=primarydark,text=white,
      minimum width=\textwidth,minimum height=1.35cm,align=center,inner sep=6pt]
  {\Large\bfseries Thème 2 — Composition des atomes et des ions\\[2pt]
   \normalsize\mdseries Niveau Moyen — Corrigé détaillé};
\end{tikzpicture}

% ---------------------------------------------------------------
\cor[un cation]
$\mathrm{Mg}^{2+}$ vient de $\nuc{24}{12}{Mg}$ ($Z=12$, $A=24$, $q=+2$).
\[ \np = Z = 12, \qquad \nz = A - Z = 24 - 12 = 12, \qquad \ee = Z - q = 12 - 2 = 10. \]
\rep $\boxed{12\ \pp,\ 12\ \nz,\ 10\ \ee}$. Le cation a \emph{perdu} $2$ électrons.

% ---------------------------------------------------------------
\cor[un anion]
$\mathrm{O}^{2-}$ vient de $\nuc{16}{8}{O}$ ($Z=8$, $A=16$, $q=-2$).
\[ \np = 8, \qquad \nz = 16 - 8 = 8, \qquad \ee = Z - q = 8 - (-2) = 8 + 2 = 10. \]
\rep $\boxed{8\ \pp,\ 8\ \nz,\ 10\ \ee}$. \\
Différence avec un cation : ici $q$ est \emph{négatif}, donc $\ee = Z - q$ \textbf{ajoute} des électrons
(l'anion en a \emph{gagné}). $\mathrm{O}^{2-}$ et $\mathrm{Mg}^{2+}$ ont d'ailleurs tous deux $10$
électrons.

% ---------------------------------------------------------------
\cor[composition d'un ion à partir de sa notation]
$\nuc{15}{7}{N}^{3-}$ : $Z = 7$, $A = 15$, $q = -3$.
\[ \np = 7, \qquad \nz = 15 - 7 = 8, \qquad \ee = 7 - (-3) = 10. \]
\rep $\boxed{7\ \pp,\ 8\ \nz,\ 10\ \ee}$.

% ---------------------------------------------------------------
\cor[un ion plus lourd]
$\nuc{132}{56}{Ba}^{2+}$ : $Z = 56$, $A = 132$, $q = +2$.
\[ \np = 56, \qquad \nz = 132 - 56 = 76, \qquad \ee = 56 - 2 = 54. \]
\rep $\boxed{56\ \pp,\ 76\ \nz,\ 54\ \ee}$.

% ---------------------------------------------------------------
\cor[atome contre ion, et le piège du noyau]
$\nuc{45}{21}{Sc}$ : $Z = 21$, $A = 45$.
\begin{enumerate}[label=\alph*)]
  \item atome neutre $\to \ee = Z = 21$ ; cation $\mathrm{Sc}^{3+} \to \ee = 21 - 3 = 18$.
  \item \textbf{Non} : le noyau est intouchable. $\np = 21$ et $\nz = 45 - 21 = 24$ dans les deux cas.
  \item \rep $\boxed{0}$ électron dans le noyau ! Le noyau ne contient \emph{que} des nucléons (protons +
        neutrons) ; les $18$ électrons de $\mathrm{Sc}^{3+}$ sont tous dans le nuage électronique.
\end{enumerate}

\begin{bilan}
Pour $\nuc{A}{Z}{X}^{q}$ : $\np = Z$, $\nz = A - Z$, $\ee = Z - q$ ($q$ signé). Former un ion ne touche
\textbf{que les électrons} ; le noyau ($Z$, $\nz$) est inchangé. Nombre d'électrons \emph{dans le noyau}
$= 0$, toujours.
\end{bilan}

\end{document}
